\documentclass[french]{book}

\usepackage[utf8]{inputenc}
\usepackage[T1]{fontenc}
\usepackage{lmodern}
\usepackage{geometry}
\usepackage{graphicx}
\usepackage{babel}
\usepackage{amsmath}
\usepackage{amsthm}
\usepackage{amssymb}
\usepackage{hyperref}
\usepackage{stmaryrd}
\usepackage{enumitem}
\usepackage{tcolorbox}
\usepackage{dsfont}
\usepackage{multirow}
\usepackage{etoc}
\usepackage{titlesec}
\usepackage{esint}
\usepackage{cancel}
\usepackage{mathdots}

\setlist[itemize]{label=$\bullet$}


\renewcommand{\P}{\mathbb{P}}
\newcommand{\N}{\mathbb{N}}
\newcommand{\Z}{\mathbb{Z}}
\newcommand{\Q}{\mathbb{Q}}
\newcommand{\R}{\mathbb{R}}
\newcommand{\C}{\mathbb{C}}
\newcommand{\K}{\mathbb{K}}
\renewcommand{\L}{\mathbb{L}}
\newcommand{\U}{\mathbb{U}}
\newcommand{\st}{^*}
\newcommand{\tq}{\middle|}
\newcommand{\inv}{^{-1}}
\newcommand{\E}{\mathbb{E}}
\newcommand{\F}{\mathbb{F}}
\newcommand{\V}{\mathbb{V}}
\renewcommand{\ker}{\text{Ker}}
\newcommand{\im}{\text{Im}}
\newcommand{\card}{\text{Card}}
\newcommand{\Tr}{\text{Tr}}

\begin{document}

\chapter{Polynômes et fractions rationnelles}

\subsection{Exercice : critère d'Eisenstein}
Soit $p$ un nombre premier et $P\in\Z[X]$ qu'on écrit sous la forme $$P=\sum_{k=0}^{n}a_kX^k$$ Supposons que l'on ait :
\begin{itemize}
	\item $\forall k\in\llbracket0,n-1\rrbracket,\ p\mid a_k$ ;
	\item $p^2\nmid a_0$.
	\item $p\nmid a_n$
\end{itemize}
Montrer que $P$ est irréductible dans $\Z[X]$.

\subsection{Exercice : polynôme à valeurs dans les nombres premiers}
On note $\mathcal{P}$ l'ensemble des nombres premiers. Trouver tous les polynômes $A \in \Z[X]$ tels que $\forall n \in \N, \ A(n) \in \mathcal{P}$.

\chapter{Révisions d'algèbre linéaire}

\chapter{Révisions sur les matrices}

\chapter{Suites et séries numériques}

\chapter{Réduction des endomorphismes}

\chapter{Espaces euclidiens}

\chapter{Topologie}

\subsection{Problème : autour de la propriété de Borel-Lebesgue}
On se donne $K$ un compact.
\begin{enumerate}
	\item  Précompacité : montrer que pour tout $\varepsilon>0$, il existe une partie finie $C$ de $K$ telle que $$K\subset\bigcup_{x\in C}B(x,\varepsilon)$$
	\item Propriété de Borel-Lebesgue :
	\begin{enumerate}[label=\alph*.]
		\item Soit $(O_i)_{i\in I}$ une famille d'ouverts de $K$ recouvrant $K$. Montrer qu'il existe $\varepsilon>0$ tel que $$\forall x\in K,\ \exists i\in I\ : \ B_o(x,\varepsilon)\subset O_i$$
		\item Montrer que pour toute famille $(O_i)_{i\in I}$ d'ouverts de $A$ recouvrant $K$, il existe une partie finie $J$ de $I$ telle que $(O_j)_{j\in J}$ recouvre $K$. On pourra utiliser la précompacité de $K$.
		\item Montrer que la propriété de la deuxième question est suffisante pour que $K$ soit compacte. On pourra utiliser le fait que l'ensemble des valeurs d'adhérence d'une suite $(u_n)$ est $\displaystyle\bigcap_{n\in\N}\overline{\{u_k\mid k\geqslant n\}}$
	\end{enumerate}
	\item Applications :
	\begin{enumerate}[label=\alph*.]
		\item Montrer que les idéaux maximaux de $\mathcal{C}^0([0,1],\R)$ sont de la forme : $$I_x=\left\{f\in\mathcal{C}^0([0,1],\R)\ \mid\ f(x)=0\right\}$$
		\item Soit $E$ un espace métrique compact et $f:E\to\R$ une fonction localement bornée, c'est-à-dire que pour tout $x\in E$, il existe un voisinage $V_x$ de $x$ tel que pour tout $y\in V_y$, on ait $\lvert f(y)\rvert\leqslant M_x$. Montrer que $f$ est bornée.
		\item On admet le théorème de Kakutani : pour tout convexe compact non vide $K$ d'un espace vectoriel normée $E$ et toute application affine continue $u$ qui stabilise $K$, $u$ possède un point fixe. Montrer qu'une famille $(u_i)_{i\in I}$ d'applications affines qui commutent deux à deux et stabilisent toutes un compact convexe non vide $K$ possèdent un point fixe commun.
	\end{enumerate}
\end{enumerate}

\subsection{Problème : autour du théorème de Baire}
Soit $A$ une partie localement compacte d'un espace vectoriel normé (\textit{ie} dans laquelle tout point possède un voisinage compact).
\begin{enumerate}
	\item Soit $(O_n)$ une suite d'ouverts de $A$, tous denses dans $A$. Montrer que $$\Omega=\bigcap_{n\in\N} O_n$$ est dense dans $A$. Est-il nécessairement ouvert ?
	\item En déduire que si une suite $(F_n)$ de fermés de $A$ a une réunion d'intérieur non vide, alors l'un des $F_n$ est d'intérieur non vide.
\end{enumerate}

\subsection{Problème : connexité par arcs}
Montrer que l'ensemble $\mathcal{GL}_n(\C)$ est connexe par arcs.

\subsection{Exercice : point fixe de Kakutani}
Soit $K$ un compact convexe non vide et $u$ une application affine continue qui stabilise $K$. Montrer que $u$ admet un point fixe.

\subsection{Exercice : sous-groupes fermés de $(\C\st,\times)$}
Déterminer les sous-groupes fermés de $(\C\st,\times)$.

\subsection{Exercice : normes sur les polynômes}
Montrer que pour tout polynôme $P\in\R[X]$, il existe une normes sur $\R[X]$ qui fait converger la suite $(X^k)_{k\in\N}$ vers $P$.

\subsection{Exercice : classe de similitude et matrices nilpotentes}
Montrer qu'une matrice $A\in\mathcal{M}_n(\K)$ est nilpotente si, et seulement si, la matrice nulle est adhérente à la classe de similitude de $A$.

\subsection{Exercice : la dérivation n'est pas continue}
Montrer qu'il n'existe pas de norme sur $\mathcal{C}^{\infty}([0,1],\R)$ pour laquelle la dérivation soit continue.

\subsection{Exercice : crochet de Lie égal à l'identité}
Soit $f$ et $g$ deux endomorphismes d'un espace vectoriel normé $E\ne\{0\}$ tels que $f\circ g-g\circ f=\text{Id}_E$.
\begin{enumerate}
	\item Calculer $f\circ g^n-g^n\circ f$ pour tout $n\in\N$ et en déduire que $f$ et $g$ ne peuvent pas être simultanément continues.
	\item Donner deux bonnes raisons de la non-existence d'un tel couple $(f,g)$ lorsque $E$ est de dimension finie.
\end{enumerate}

\subsection{Exercice : point fixe d'une application contractante}
Soit $X$ une partie fermée non vide d'un espace vectoriel de dimension finie, $k\in[0,1[$ et $f:X\to X$ une application $k$-lipschitzienne. Montrer que $f$ possède un unique point fixe. Indication : on pourra, étant donné $x\in X$, considérer la série $$\sum (f^{n+1}(x)-f^n(x))$$

\subsection{Exercice : un point fixe}
Soit $C$ une partie convexe compacte non vide et $f:C\to C$ 1-lipschitzienne. Montrer que $f$ admet au moins un point fixe. On pourra commencer par le cas où $f$ est $k$-lipschitzienne avec $k<1$.

\subsection{Exercice : vers la connexité}
Soit $A$ une partie non triviale d'un espace vectoriel normé $E$. Montrer que la frontière de $A$ est non vide.

\end{document}